\documentclass[../dayan_api_reference.tex]{subfiles}

\begin{document}

\section{dapi API}\label{sec:api-dapi}

dapi 是 kernel 和 Device 子 Proc 可直接调用的稳定接口。公开面由
\code{device/include/api} 下的头文件定义，动态链接允许的函数集合由
\code{device/dayan-api.symbols} 固定。本章逐项记录这 61 个可链接符号，并列出公开
类型、常量与编译属性；BL1 scheduler、Proc 和资源实现细节不属于 dapi。

\subsection{头文件与编译属性}

\begin{dayantable}{dapi 公开头文件}{tab:api-dapi-headers}{X[34,l] X[66,l]}
  头文件 & 内容 \\
  \code{dayan_sys.h} & Proc、launch 坐标、URAM、Cluster barrier、执行单元等待与资源 API。 \\
  \code{dayan_attrs.h} & dapi / Device libc 的 visibility、noreturn 和 restrict 属性。 \\
  \code{stdio.h} / \code{stdlib.h} / \code{string.h} / \code{math.h} & Device libc 子集。 \\
  \code{stdint.h} / \code{stddef.h} / \code{stdbool.h} / \code{stdarg.h} / \code{limits.h} & freestanding 基础类型和宏。 \\
\end{dayantable}

\begin{dayantable}{dapi 编译属性}{tab:api-dapi-attrs}{X[36,l] X[64,l]}
  宏 & 含义 \\
  \code{DAYAN_API} & 将 Dayan 系统函数放入公开 dapi section 并标记为 used。 \\
  \code{DAYAN_LIBC_API} & 将 Device libc 函数放入公开 libc section 并标记为 used。 \\
  \code{DAYAN_NORETURN} & 声明函数不返回。 \\
  \code{DAYAN_RESTRICT} & C/C++ 兼容的 restrict 修饰。 \\
\end{dayantable}

\subsection{Kernel 环境与资源描述}

\begin{dayantable}{Kernel 环境 API}{tab:api-dapi-environment}{X[43,l] X[57,l]}
  API & 含义 \\
  \code{shared} & 绑定 \code{gp} 的 Host-owned coherent allocation 基址；kernel 只能读取该指针本身，可在所指区间放跨 Hart 状态。 \\
  \code{DAYAN_KERNEL_RESOURCE_DESC_VERSION} & 当前静态 kernel 资源描述版本 1。 \\
  \code{DAYAN_HART_MAX_PROCS} & 每个 Hart 最多 16 个用户 Proc slot。 \\
  \code{DAYAN_KERNEL_MAX_CHILD_PROCS} & root Proc 最多声明 15 个子 Proc。 \\
  \code{DAYAN_PROC_NAME_SIZE} & Proc 名称存储为 16 B。 \\
  \code{DAYAN_URAM_ALLOCATION_GRANULE} & URAM 资源以 1024 B 对齐和分配。 \\
  \code{dayan_kernel_resource_desc_t} & dycc/Runtime 消费的版本化常量描述：最大子 Proc 数、Atomic Cache 字节数及 URAM 字节数。 \\
  \code{dayan_proc_entry_t} & 子 Proc 入口类型 \code{void (*)(void *)}。 \\
  \code{dayan_cluster_barrier_t} & 四 Core barrier 的 16-B、零初始化 shared storage。 \\
  \code{dayan_resource_state} & \code{DAYAN_RESOURCE_FREE}、\code{DAYAN_RESOURCE_BUSY}、\code{DAYAN_RESOURCE_READY} 三态资源状态机。 \\
  \code{dayan_task_axis} & \code{DAYAN_TASK_X}、\code{DAYAN_TASK_Y}、\code{DAYAN_TASK_Z} 三个 logical task 轴。 \\
  \code{dayan_resource_t} & 同一 Hart 内的资源对象；字段私有，用户只负责存储生命周期。 \\
  \code{dayan_shared_resource_t} & 跨 Hart 资源对象；必须位于 \code{shared} 指向的分配内。 \\
  \code{list_head} & \code{dayan_resource_t} 的公开布局依赖类型；不作为独立容器 API 使用。 \\
\end{dayantable}

Kernel 使用 \code{DAYAN_KERNEL_DESC}（见 Runtime 章节）生成一个
\code{dayan_kernel_resource_desc_t} 常量。Atomic Cache/URAM 数量在 Host launch 前一次性
分配，子 Proc slot 在 root Proc 建立时一次性预留。

\subsection{Proc 生命周期}

\begin{dayantable}{Proc dapi}{tab:api-dapi-proc}{X[39,l] X[61,l]}
  API & 行为 \\
  \code{dayan_yield} & 当前 Proc 主动让出调度机会。 \\
  \code{dayan_create_proc} & 在同一物理 Core 的预留 slot 中创建具名直接子 Proc；成功返回正 Proc ID。名称必须非空、NUL 结尾并能放入 16 B。 \\
  \code{dayan_join} & 等待指定直接子 Proc，按需返回其 64-bit result，并回收 zombie。 \\
  \code{dayan_join_all} & 等待并回收当前 Proc 的全部直接子 Proc。 \\
  \code{dayan_proc_exit} & 以 64-bit result 结束当前 Proc，不返回。 \\
\end{dayantable}

父 Proc 退出前必须等待其子 Proc；root Proc 退出时 BL1 将结果通知 Host。子 Proc 的创建、
退出和 slot 回收完全在 Device 内完成，不产生 Host create/exit RPC。

\subsection{Launch、URAM 与 Cluster 同步}

\begin{dayantable}{执行环境查询与同步}{tab:api-dapi-launch}{X[39,l] X[61,l]}
  API & 行为 \\
  \code{dayan_uram_base} & 返回当前 resident kernel instance 独占 URAM 区间的基址；未处于 Proc 时返回空。 \\
  \code{dayan_uram_size} & 返回该 URAM 区间字节数。 \\
  \code{dayan_task_id} & 返回 X/Y/Z 轴的当前 logical task ID；非法轴返回 0。 \\
  \code{dayan_task_dim} & 返回 X/Y/Z 轴的 launch grid 大小；非法轴返回 0。 \\
  \code{dayan_task_mode} & 返回 1（单 Core）或 4（Cluster 协作）。 \\
  \code{dayan_cluster_barrier} & 对 shared allocation 中的 barrier 执行可复用四 Core rendezvous；仅 mode 4 且 participants=4 合法。 \\
\end{dayantable}

logical task API 不暴露物理 Cluster/Core 编号。四 Core mode 下 Runtime 要求
\code{task_dim.x == 4}，X ID 对应同一 Cluster 内四个参与者。

\subsection{执行单元完成等待}

\begin{dayantable}{执行单元等待 dapi}{tab:api-dapi-wait}{X[37,l] X[63,l]}
  API & 行为 \\
  \code{dayan_wait_lsu} & 等待当前 Proc 此前提交的 LSU 工作完成，并处理 scheduler produce token。 \\
  \code{dayan_wait_dma} & LSU/DMA 等待的兼容名称，当前与 \code{dayan_wait_lsu} 等价。 \\
  \code{dayan_wait_mat} & 等待此前提交的矩阵单元工作完成。 \\
  \code{dayan_wait_vec} & 等待此前提交的向量单元工作完成。 \\
  \code{dayan_wait_macro} & 等待此前提交的 macro 单元工作完成。 \\
\end{dayantable}

这些函数替代 kernel 内手写 barrier 宏。它们只保证当前执行单元此前工作的完成顺序；
Core 间发布/消费关系应使用 shared resource 或 Cluster barrier 表达。

\subsection{单 Hart 资源}

\code{dayan_resource_t} 供同一 Hart 上的 producer/consumer Proc 协作。资源循环为
FREE $\rightarrow$ BUSY $\rightarrow$ READY $\rightarrow$ FREE；当 producer count 为
零时 alloc 直接进入 READY。所有函数成功返回 0，失败返回负 errno。

\begin{dayantable}{单 Hart resource dapi}{tab:api-dapi-resource}{X[44,l] X[56,l]}
  API & 行为 \\
  \code{dayan_resource_init} & 初始化对象为 FREE，并建立本 Hart 的等待队列。 \\
  \code{dayan_resource_alloc} & 仅从 FREE 开始一代资源，设置 producer/consumer count；consumer count 必须非零。 \\
  \code{dayan_resource_produce} & 递减 producer count；归零时发布 READY 并唤醒 waiter。 \\
  \code{dayan_resource_consume} & 递减 consumer count；归零时释放为 FREE 并唤醒 waiter。 \\
  \code{dayan_resource_wait_ready} & 阻塞当前 Proc，直到资源达到 READY。 \\
  \code{dayan_resource_wait_free} & 阻塞当前 Proc，直到资源回到 FREE。 \\
\end{dayantable}

\subsection{跨 Hart 资源}

\code{dayan_shared_resource_t} 表达 Cluster 内跨 Hart 的发布/消费。对象必须位于当前 kernel
的 \code{shared} allocation 且满足自然对齐；其状态更新使用 Atomic Cache 可见的原子
操作，等待侧以 scheduler poll 让出 Core。producer/consumer 剩余计数打包在同一个
64-bit 状态字中，分配使用 \code{amocas.d}，发布和消费使用 \code{amoadd.d}，不使用
自旋锁。

\begin{dayantable}{Shared resource dapi}{tab:api-dapi-shared-resource}{X[47,l] X[53,l]}
  API & 行为 \\
  \code{dayan_shared_resource_init} & 校验 shared 范围，并将跨 Hart 对象初始化为 FREE。 \\
  \code{dayan_shared_resource_alloc} & 使用 \code{amocas.d} 从 FREE 开始一代资源并设置计数。 \\
  \code{dayan_shared_resource_produce} & 使用 \code{amoadd.d} 原子递减 producer count，最后一个 producer 发布 READY。 \\
  \code{dayan_shared_resource_consume} & 使用 \code{amoadd.d} 原子递减 consumer count，最后一个 consumer 发布 FREE。 \\
  \code{dayan_shared_resource_wait_ready} & 可调度地轮询，直到跨 Hart 资源为 READY。 \\
  \code{dayan_shared_resource_wait_free} & 可调度地轮询，直到跨 Hart 资源为 FREE。 \\
\end{dayantable}

\subsection{Device stdio}

\begin{dayantable}{Device stdio API}{tab:api-dapi-stdio}{X[35,l] X[65,l]}
  API & 行为 \\
  \code{FILE} & 不透明 stream 类型；当前公开 stream 只有 \code{stdout}。 \\
  \code{stdout} & 当前 Proc 的 Host 输出缓冲。 \\
  \code{EOF} & stdio 失败返回值 $-1$。 \\
  \code{putchar} & 向当前 Proc 输出一个字节；失败返回 EOF。 \\
  \code{puts} & 输出字符串及换行；失败返回 EOF。 \\
  \code{fflush} & 将当前 Proc 的累积输出通过 RPC 发给 Host；参数可为空或 stdout。 \\
  \code{printf} & 格式化写入当前 Proc 输出缓冲。 \\
  \code{vprintf} & \code{va_list} 版本的 printf。 \\
  \code{snprintf} & 有界格式化到字符数组，并保证 size 非零时结尾 NUL。 \\
  \code{vsnprintf} & \code{va_list} 版本的 snprintf。 \\
\end{dayantable}

Proc 的 printf buffer 有固定容量，写满时自动 flush；Proc 结束时剩余内容也会 flush，Host
随后清空该 Proc 的输出状态。

\subsection{Device 动态内存}

\begin{dayantable}{Device stdlib API}{tab:api-dapi-stdlib}{X[34,l] X[66,l]}
  API & 行为 \\
  \code{malloc} & 从 BL1 管理的当前 Hart heap 分配至少指定字节数。 \\
  \code{calloc} & 检查乘法溢出后分配并清零 count $\times$ size 字节。 \\
  \code{realloc} & 调整已有分配；返回值可能与原地址不同。 \\
  \code{free} & 释放分配；空指针可安全传入。 \\
  \code{abort} & 停止当前 Hart，不返回。 \\
  \code{exit} & 当前实现忽略 status 并停止 Hart，不返回；正常 Proc 返回应使用 \code{dayan_proc_exit}。 \\
\end{dayantable}

\subsection{Device 字符串与内存}

\begin{dayantable}{Device string API}{tab:api-dapi-string}{X[34,l] X[66,l]}
  API & 行为 \\
  \code{memcpy} & 复制不重叠的 size 个字节，返回 dst。 \\
  \code{memmove} & 复制可能重叠的 size 个字节，返回 dst。 \\
  \code{memset} & 用 value 的低 8 bit 填充 size 个字节，返回 dst。 \\
  \code{memcmp} & 按 unsigned byte 字典序比较两个内存区。 \\
  \code{strlen} & 返回 NUL 结尾字符串的长度。 \\
  \code{strnlen} & 返回字符串长度，但最多检查 maxlen 字节。 \\
  \code{strcpy} & 复制包含结尾 NUL 的字符串，返回 dst。 \\
  \code{strncpy} & 最多复制 size 字节并按标准 strncpy 规则补零。 \\
  \code{strcmp} & 按 unsigned char 比较两个 NUL 结尾字符串。 \\
  \code{strncmp} & 最多比较 size 个字符。 \\
  \code{strchr} & 返回字符第一次出现的位置或空指针。 \\
  \code{strrchr} & 返回字符最后一次出现的位置或空指针。 \\
\end{dayantable}

\subsection{Device 数学函数}

\begin{dayantable}{Device math API}{tab:api-dapi-math}{X[34,l] X[66,l]}
  API & 行为 \\
  \code{fabsf} & 单精度绝对值。 \\
  \code{copysignf} & 取 value 的幅值与 sign 的符号位。 \\
  \code{fminf} & 单精度最小值，按实现处理 NaN。 \\
  \code{fmaxf} & 单精度最大值，按实现处理 NaN。 \\
  \code{sqrtf} & 单精度平方根。 \\
  \code{rsqrtf} & 单精度平方根倒数。 \\
  \code{expf} & 单精度自然指数。 \\
  \code{logf} & 单精度自然对数。 \\
  \code{INFINITY} / \code{NAN} / \code{HUGE_VALF} & 单精度特殊值常量。 \\
  \code{isnan} / \code{isinf} / \code{isfinite} / \code{signbit} & 编译器 builtin 分类宏。 \\
\end{dayantable}

\subsection{Freestanding 基础定义}

这些头文件是面向 RISC-V Device 的 freestanding 子集，不承诺 Host libc 的完整功能、
locale、文件系统、errno 或线程安全语义。

\begin{dayantable}{基础类型与宏}{tab:api-dapi-foundation}{X[43,l] X[57,l]}
  头文件 & 对外定义 \\
  \code{stdint.h} types & \code{int8_t}、\code{uint8_t}、\code{int16_t}、\code{uint16_t}、\code{int32_t}、\code{uint32_t}、\code{int64_t}、\code{uint64_t}、\code{intptr_t}、\code{uintptr_t}、\code{intmax_t} 和 \code{uintmax_t}。 \\
  \code{stdint.h} 8/16-bit limits & \code{INT8_MIN}、\code{INT8_MAX}、\code{UINT8_MAX}、\code{INT16_MIN}、\code{INT16_MAX}、\code{UINT16_MAX}。 \\
  \code{stdint.h} 32/64-bit limits & \code{INT32_MIN}、\code{INT32_MAX}、\code{UINT32_MAX}、\code{INT64_MIN}、\code{INT64_MAX}、\code{UINT64_MAX}。 \\
  \code{stdint.h} pointer limits & \code{INTPTR_MAX}、\code{UINTPTR_MAX} 和 \code{SIZE_MAX}。 \\
  \code{stddef.h} & \code{size_t}、\code{ptrdiff_t}、\code{wchar_t}、\code{max_align_t}、\code{NULL} 和 \code{offsetof}。 \\
  \code{stdbool.h} & C 模式下的 \code{bool}、\code{true}、\code{false}。 \\
  \code{stdarg.h} & \code{va_list}、\code{va_start}、\code{va_arg}、\code{va_end} 和 \code{va_copy}。 \\
  \code{limits.h} & \code{CHAR_BIT}、\code{SCHAR_MIN}、\code{SCHAR_MAX}、\code{UCHAR_MAX}、\code{SHRT_MIN}、\code{SHRT_MAX}、\code{USHRT_MAX}、\code{INT_MIN}、\code{INT_MAX}、\code{UINT_MAX}、\code{LONG_MIN}、\code{LONG_MAX} 和 \code{ULONG_MAX}。 \\
\end{dayantable}

\end{document}
